\chapter{검증 결과(10): 최종 설계로 다시 계산한 결론}

6–12장과 15장의 결론 상당수는 150 kHz Stage A(위치당 6 mΩ 가정, 출력용량 손실 4·E\textsubscript{oss}·f\textsubscript{s}, C\textsubscript{A} 400 µF)로 계산했으므로 A25R로 다시 계산했다(E36, TB28). 외부 검토 뒤의 재계산(E38)도 이 장에 모은다. 각 장의 본문은 최종 값을 쓰고, 바뀐 값은 표 \ref{t:recalc}에 모았다.

\section{다시 계산한 범위와 방법}

A25R Stage A 모델은 스위칭 모의 격자(640–920 V × 위상당 1–50 A, 충전 375 A)의 손실·V\textsubscript{A} 강하 표, 위상 셰딩(멈춤 20 A), 모듈 출력저항 R\textsubscript{A} 9.0 mΩ(150 kHz 설계 10.2 mΩ), C\textsubscript{A} 2.4 mF이며, 판정 기준은 원 시험의 것을 실행 전에 그대로 고정했다. 범위는 E15, E25, E18, E16, E23, 클램프 없는 충전 중 정지 8사건(E26), 이중 고장 4사례(E29), TB28이고, 충전기 급정지(정격 충전 중 10 µs에 0 A, 640·800·920 V × 공칭·표본, 기준: 위상 첨두 ≤ 1.5 × P\textsubscript{rated}, V\textsubscript{A} ±10 V·≤ 470 V)를 새로 더했다.


이어서 외부 검토(2026-10-29)의 지적으로 세 가지를 바꿔 다시 계산했다(E38). 손실 비교의 Stage C를 최종 설계와 같은 규칙의 고정이득 5 kW 직류변압기로 바꿨고, 충전 손실을 실제 위상 전류의 스위칭 모의로 직접 계산했으며, “150 kW 충전 손실”을 “375 A 정전류(입력 120·150·172.5 kW)”로 고쳤다. 640 V에서 150 kW를 내려면 469 A가 필요해 충전 정격(375 A) 밖이기 때문이다.

\Needspace{12\baselineskip}\section{150 kHz 값과 최종 값의 대조}


{\footnotesize
\begin{longtable}{L{\dimexpr 0.24\linewidth-2\tabcolsep\relax}L{\dimexpr 0.22\linewidth-2\tabcolsep\relax}L{\dimexpr 0.22\linewidth-2\tabcolsep\relax}L{\dimexpr 0.25\linewidth-2\tabcolsep\relax}L{\dimexpr 0.07\linewidth-2\tabcolsep\relax}}
\caption{150 kHz 설계 값, A25R 재계산 값(E36), 최종 값}\label{t:recalc}\\
\toprule
\tabhead \textbf{항목} & \textbf{150 kHz(중간보고서)} & \textbf{A25R(E36)} & \textbf{최종(E38–E44)} & \textbf{장} \\
\midrule
\endfirsthead
\toprule
\tabhead \textbf{항목} & \textbf{150 kHz(중간보고서)} & \textbf{A25R(E36)} & \textbf{최종(E38–E44)} & \textbf{장} \\
\midrule
\endhead
\bottomrule
\endlastfoot
④ 보조 경로 손실(주행 대표) & 85.0 W & 94.4 W & 101.5 W(최종 Stage C) & 6 \\
① / ② / ③ / ③′ & 114.6 / 123.8 / 133.8 / 114.1 W & 114.6 / 133.3 / 143.2 / 123.4 W & 112.4 / 134.6 / 144.5 / 130.4 W(① 소자 조건 정정 E44, 정정 전 115.9) & 6 \\
손실 순위 & ④ < ③′ < ① < ② < ③ & ④ < ① < ③′ < ② < ③ & 같음 & 6 \\
④ 스위치 VA·위치 & 953 kVA·44 & 949 kVA·44 & 1,009 kVA·42 & 6 \\
③ = ④ 손익분기(Stage B 고정 손실 배율) & 0.228 & 0.228 & 0.32 & 6 \\
375 A 충전 손실(640·800·920 V) & 0.50 kW(“150 kW”, 전압 무관 497 W) & 520·586·639 W & 같음(375 A 고정 전류의 손실 비교, 반전압 기준 120·150·172.5 kW) & 6, 15 \\
사양 팩 150 kW 충전(가치 지도) & 485 W & 566 W(제곱 환산) & 582 W(직접 계산) & 7 \\
④ 대 ⑤ 보조 손실 & 26 \% 작음 & 18 \% 작음 & 12 \% 작음(E38) → 같은 수준, 0.5 \%(E41) → −1.4 \%(E44, 명목 순서 바뀜, 같은 수준 유지) & 15 \\
④ 대 ⑥ 보조 손실 & 15 W 큼 & 24 W 큼 & 19 W 큼(E38) → 4–30 W(E41) → 4–33 W 큼(E43) & 15 \\
충전 손실 대 모터 승압 하한 & 모든 전압에서 작음 & 920 V에서 6 W 큼 & 920 V에서 12 W 큼(E44) & 15 \\
충전 손실 대 전용 승압기 & 28 \%(50 kHz) & 33 \%(50 kHz) & 0.32–0.51배(16–50 kHz, E41) & 6, 15 \\
가치 지도 B\textsubscript{420} 최선 & 200·216s·사양 팩 R2-W410 & 200s만 R2-W410, 216s·사양 팩 ①(5.8·9.1 W) & 공정 기준(E42): 216·224s·사양 팩 ① 최선, 200s ①·R2-W410 같은 수준 & 7 \\
240s R2-W460 대 ② & 15.5 W(13 \%) 작음 & 15.5 W(11 \%) 작음, VA도 작음 & 8.3 W(6 \%) 작음, VA 22 kVA 큼. 민감도에서 같은 수준(E42) & 7 \\
가치 지도 B\textsubscript{470} 최선 & 180–216s·사양 팩 ④, 224·240s R2-W460 & 같음 & 공정 기준(E42): 모든 칸 같은 수준(표 \ref{t:e18}과 각주 b; E44 정정 뒤 200·216s·사양 팩의 명목 최소 ④ → ①) & 7 \\
Stage A 공진 L 저장에너지(8위상) & 3.2 mJ(승압기의 약 1/1000) & 18.9 mJ(약 1/175) & 같음 & 7 \\
클램프 없는 470 V 도달 시간 & 0.01–1.1 ms & 0.06–2.08 ms & 같음 & 12 \\
\end{longtable}
}
\srcnote{주: 150 kHz 값의 손실은 출력용량 손실을 한 자릿수 작게 잡은 모델의 값으로, 손실 증가는 모델 정정의 결과다. 충전 램프·셰딩 문턱·클램프·이중 고장·소프트 투입의 변화는 표 \ref{t:changes}, \ref{t:a25r_mgr}, \ref{t:clamp_hist}, \ref{t:fails}과 9.6절에 있다. 자료: 검증 근거 E15, E18, E25, E26, E29, E36–E41.}

\section{바뀐 결론}

A25R 재계산(E36)에서 바뀐 결론은 세 가지였다.

\begin{enumerate}[label=(\arabic*)]
\item \textbf{충전 손실이 모터 승압 하한보다 작다는 문장(E25).} 920 V에서 639 W로 모터 승압 하한(상저항 5 mΩ)보다 커졌다(정정 전 하한 633 W 대비 6 W, 비교 소자 E\textsubscript{oss} 조건 정정 뒤 하한 627 W 대비 12 W, E44). “모든 전압에서 하한보다 작다”는 문장은 쓰지 않고, 기준점(상저항 10 mΩ)보다 모든 전압에서 작다는 것과 격자의 모든 모터 승압 조합보다 Stage A가 작아지는 상저항(640·800·920 V에서 4.0·4.7·5.2 mΩ 이상, E44 정정 값; 정정 전 3.9·4.6·5.1 mΩ)만 쓴다(15장).
\item \textbf{가치 지도의 B\textsubscript{420} 두 칸(E18).} 216s와 사양 팩의 최선 구조가 R2-W410에서 ①로 바뀌었다. R2-W410은 Stage A 손실 증가(약 9 W)를 그대로 받고, ①은 Stage A가 없어 그대로이기 때문이다. 성립 여부가 바뀐 칸은 없다.
\item \textbf{이중 고장 F2(E29).} C\textsubscript{A}가 6배라 0/4가 되었다. 방전 저항(E37)으로는 2/4였고, 최종 대책은 400 V 부하 정지다(수치는 9.6절).
\end{enumerate}

그 뒤 E38(Stage C 모델 교체, 충전 손실 직접 계산)에서는 성립·최선 칸과 결론 문장이 바뀌지 않았고, 240s R2-W460 대 ②의 문장 하나(“손실·VA 모두 작다” → “손실만 8.3 W 작고 VA는 22 kVA 크다”)만 바뀌었다. E41·E42로 바뀐 비교는 표 \ref{t:recalc} 및 7장·15장에 정리했다.


\textbf{바뀌지 않은 결론.} H5 기각(6.2절), 충전 중 보조버스(H3)와 R2 버스·48 V(8.2절, 10.4절. E16은 Stage A 정적 등가로 다시 계산했고, 그 보수성은 증명하지 않았다, 18.4절), 충전 중 정지에 클램프가 필요하다는 결론(920 V 충전에서는 클램프 없이 0.1 ms 안에 470 V에 닿는다, 12.4절. 새로 넣은 충전기 급정지(10 µs) 사례도 6/6 통과), Stage A 충전 경로의 이점(15장)은 바뀌지 않았다. H3의 조건은 사전 기동과 600 µs 램프로 바뀌었다.
